% tableview-diffable.tex — UITableView data source/delegate, cell reuse, UICollectionView,
% diffable data source + snapshots, reconfigure vs reload, CellRegistration.
% Sources (checked 2026-10-05):
%   developer.apple.com/documentation/uikit/uitableviewdatasource · .../uitableviewdelegate ·
%     .../uitableview/dequeuereusablecell(withidentifier:for:) · .../uitableviewcell/prepareforreuse()
%   developer.apple.com/documentation/uikit/uitableviewdiffabledatasource (iOS 13) ·
%     .../nsdiffabledatasourcesnapshot (reconfigureItems iOS 15, reloadItems iOS 13; identifiers
%     must be unique) · .../uicollectionview/cellregistration (iOS 14)
%   developer.apple.com/documentation/uikit/uitableviewcell/cellstyle
%   developer.apple.com/videos/play/wwdc2021/10252 "Make blazing fast lists and collection views"
% @labels: area=ios-swift kind=api platform=ios topic=ui discussion=71
% @tags: uitableview, uicollectionview, cell-reuse, dequeuereusablecell, prepareforreuse, diffable-data-source, nsdiffabledatasourcesnapshot, reconfigureitems, reloaditems, cellregistration, performbatchupdates, compositional-layout
\documentclass[8pt]{extarticle}
\usepackage{printup-sheet}

\lstdefinelanguage{SwiftSheet}{
  morekeywords={func,let,var,struct,class,final,enum,protocol,extension,return,if,else,
    guard,try,await,async,case,switch,self,some,any,where,init,in,private,override},
  sensitive=true, morecomment=[l]{//}, morestring=[b]"}

\begin{document}

\sheettitle{Table views, cell reuse \& diffable data source}{ios-swift · folio}

\oneliner{A list keeps only the \textbf{visible} cells alive and \textbf{recycles} them — so a
dequeued cell is \emph{used}, never blank. A \textbf{diffable data source} replaces index math:
you describe the \emph{whole desired state} as a \textbf{snapshot of unique \texttt{Hashable} IDs},
\texttt{apply} it, and UIKit computes the inserts/deletes/moves.}

\begin{multicols}{2}
\raggedright

\section{How it works}
\begin{itemize}
  \item \textbf{\texttt{UITableViewDataSource}} (\emph{what} to show) — required:
        \texttt{tableView(\_:numberOfRowsInSection:)}, \texttt{tableView(\_:cellForRowAt:)};
        optional \texttt{numberOfSections(in:)}. \textbf{\texttt{UITableViewDelegate}}
        (\emph{behaviour}) — \texttt{didSelectRowAt}, \texttt{heightForRowAt}, \texttt{willDisplay}.
        Both held \texttt{weak}.
  \item \textbf{Reuse}: \texttt{register(\_:forCellReuseIdentifier:)} once, then
        \texttt{dequeueReusableCell(withIdentifier:for:)} — always returns a cell (crashes if
        unregistered). Old \texttt{dequeueReusableCell(withIdentifier:)} may return \texttt{nil}.
        A cell leaving the screen goes to a per-identifier \textbf{reuse queue}; \sloppy
        \texttt{prepareForReuse()} runs before it comes back.
  \item \textbf{Stale content} = a recycled cell still shows the last row's image/check/colour,
        because \texttt{cellForRowAt} set it only on \emph{some} paths, or an async image load
        finished \emph{after} reuse. Fix: set \textbf{every} property on every path; in
        \texttt{prepareForReuse} reset transient state + \textbf{cancel} the load; on completion
        check the cell still shows the same ID.
  \item \textbf{Manual updates}: change the model \emph{first}, then \texttt{performBatchUpdates}
        with \texttt{insertRows}/\texttt{deleteRows}. Counts must add up or
        \texttt{NSInternalInconsistencyException} (``invalid number of rows'').
  \item \textbf{\texttt{UICollectionView}} = same data source/delegate/reuse, plus a \textbf{layout}
        object: \texttt{UICollectionViewFlowLayout} or \texttt{UICollectionViewCompositionalLayout}
        (item $\to$ group $\to$ section); lists via \texttt{UICollectionLayoutListConfiguration} (iOS 14).
  \item \textbf{Diffable} (iOS 13): \texttt{UITableViewDiffableDataSource} /
        \texttt{UICollectionViewDiffableDataSource<SectionID, ItemID>}; both generic params
        \textbf{\texttt{Hashable}}; created with a \texttt{cellProvider} closure.
        \texttt{NSDiffableDataSourceSnapshot} is a \emph{struct}: \texttt{appendSections},
        \texttt{appendItems(\_:toSection:)}, \texttt{deleteItems}, \texttt{moveItem(\_:beforeItem:)},
        then \texttt{apply(\_:animatingDifferences:)}. Edit the live one via
        \texttt{dataSource.snapshot()}. Map positions with \texttt{itemIdentifier(for:)} /
        \texttt{indexPath(for:)}.
  \item \textbf{Same ID, new content}: \texttt{reconfigureItems} (iOS 15) re-runs the config on the
        \emph{existing} cell — cheap, keeps state. \texttt{reloadItems} (iOS 13) throws the cell away
        and dequeues a new one. Prefer reconfigure unless the cell \emph{type} changes.
  \item \textbf{\texttt{CellRegistration}} (iOS 14): typed cell + config closure, no string IDs;
        \texttt{dequeueConfiguredReusableCell(using:for:item:)}. Headers need
        \texttt{supplementaryViewProvider} (collection) or a subclass overriding
        \texttt{titleForHeaderInSection} (table).
\end{itemize}

\section{Example}
\begin{lstlisting}[language=SwiftSheet]
enum Section: Hashable { case main }
let reg = UICollectionView.CellRegistration
    <UICollectionViewListCell, Todo.ID> { cell, _, id in
  var c = cell.defaultContentConfiguration()
  c.text = store[id]?.title        // read LIVE data by id
  cell.contentConfiguration = c    // set EVERY property
}
ds = UICollectionViewDiffableDataSource<Section, Todo.ID>(
  collectionView: cv) { cv, ip, id in
  cv.dequeueConfiguredReusableCell(using: reg, for: ip, item: id) }
var snap = NSDiffableDataSourceSnapshot<Section, Todo.ID>()
snap.appendSections([.main])
snap.appendItems(todos.map(\.id))  // ids UNIQUE
ds.apply(snap, animatingDifferences: true)
\end{lstlisting}

\columnbreak

\section{The reuse pool}
\begin{tikzpicture}[sheet]
  \tikzset{row/.style={cell, minimum width=24mm, minimum height=4.6mm, font=\ttfamily\scriptsize}}
  \node[draw=sheetBlue, thick, rounded corners=3pt, minimum width=28mm, minimum height=28.5mm,
        fill=sheetBlue!4] (scr) at (0,0) {};
  \node[note, font=\scriptsize\itshape] at (0,1.2) {screen: visible rows only};
  \node[row, fill=black!4, text=sheetGrey] (r6) at (0,2.0) {row 6 (scrolled off)};
  \node[row] at (0,0.8) {row 7};
  \node[row] at (0,0.27) {row 8};
  \node[row] at (0,-0.26) {row 9};
  \node[row] at (0,-0.79) {row 10};
  \node[row, draw=sheetGreen, fill=sheetGreen!10] (r11) at (0,-2.0) {row 11 (appearing)};
  \node[box, fill=sheetOrange!8, draw=sheetOrange, text width=27mm, align=left, font=\footnotesize]
       (pool) at (4.5,0) {\textbf{reuse queue}\\[-1pt]per identifier\\[1pt]
       \texttt{prepareForReuse()}\\[-1pt]{\scriptsize reset + cancel loads}};
  \draw[hot] (r6.east) to[out=0,in=90] node[above right, note, pos=0.3]{enqueue} (pool.north);
  \draw[hot] (pool.south) to[out=-90,in=0] (r11.east);
  \node[note, align=left, anchor=west] at (4.1,-1.6)
       {\texttt{dequeueReusableCell}\\[-1pt]\texttt{(withIdentifier:for:)}};
  \node[note, text=sheetRed, anchor=north, align=center, text width=75mm] at (2.3,-2.3)
       {same object, still holds row 6's image $\to$ \texttt{cellForRowAt} sets \textbf{all}};
\end{tikzpicture}

\section{Snapshot diff: IDs, not content}
\begin{tikzpicture}[sheet]
  \tikzset{sid/.style={cell, minimum width=12mm, minimum height=4.8mm}}
  \node[note] at (0,0.5) {old snapshot};
  \node[note] at (3.0,0.5) {new snapshot};
  \node[sid] (a0) at (0,0) {A};   \node[sid] (a1) at (3.0,0) {A};
  \node[sid, draw=sheetRed, fill=sheetRed!8] (b0) at (0,-0.6) {B};
  \node[sid, fill=sheetBlue!8] (c0) at (0,-1.2) {C};
  \node[sid, fill=sheetBlue!8] (c1) at (3.0,-0.6) {C\,\textsuperscript{*}};
  \node[sid, draw=sheetGreen, fill=sheetGreen!10] (d1) at (3.0,-1.2) {D};
  \draw[flow] (a0) -- node[above, note]{kept} (a1);
  \draw[flow, sheetBlue] (c0.east) -- node[below, sloped, note, pos=0.4]{move} (c1.west);
  \node[note, text=sheetRed, left=1mm of b0] {delete};
  \node[note, text=sheetGreen, right=1mm of d1] {insert};
  \node[note, align=left, text width=33mm, anchor=north west] at (3.75,0.3)
       {\textsuperscript{*}title changed, id same $\to$ the diff sees \emph{nothing}.
        \texttt{reconfigureItems([C])}, then \texttt{apply}.};
  \node[note, align=left, text width=76mm, anchor=north west, text=sheetRed] at (-1.0,-1.65)
       {Whole model \emph{as} the id (\texttt{Hashable} over all fields) $\to$ an edit changes the
        hash $\to$ delete + insert: flicker, lost selection. Two equal ids $\to$ crash.};
\end{tikzpicture}

\section{Traps}
\begin{itemize}
  \trap{\textbf{``Diffable detects content changes.''} No — it diffs \textbf{identifiers} only.
        Same ID + new title = no update until \texttt{reconfigureItems}/\texttt{reloadItems}.}
  \trap{\textbf{Duplicate IDs} anywhere in the snapshot (across sections too) $\to$ runtime crash.
        Never use a display name as the ID.}
  \trap{Caching an \texttt{IndexPath} across an \texttt{apply} — positions shift; convert to the ID
        with \texttt{itemIdentifier(for:)} immediately.}
  \trap{Config closure capturing a model copy $\to$ reconfigure shows old data. Look it up by ID.}
  \trap{Apply from main \emph{or} always the same background queue — never mix.}
  \trap{\texttt{prepareForReuse} is not the place to set content — that is \texttt{cellForRowAt}.}
\end{itemize}

\section{\texttt{UITableViewCell.CellStyle}}
{\footnotesize
\begin{tabular}{@{}l l@{}}
\toprule
\texttt{.default} & title only (+ image) \\
\texttt{.subtitle} & title, smaller grey subtitle \emph{below} \\
\texttt{.value1} & title left, detail right-aligned (Settings) \\
\texttt{.value2} & small blue label, then detail (Contacts) \\
\bottomrule
\end{tabular}}\pdfonly{\\[1pt]}
{\scriptsize iOS 14+: \texttt{defaultContentConfiguration()} replaces \texttt{textLabel} /
\texttt{detailTextLabel}.}

\section{Remember}
\textbf{``The pool gives you a \emph{used} cell — set everything, assume nothing.''}
\textbf{``Snapshot = IDs; content = your store; reconfigure = repaint.''}


\end{multicols}

\noindent{\footnotesize\color{sheetGrey}\textit{Related:} UICollectionView (compositional layout) · Auto Layout (self-sizing cells:
\texttt{automaticDimension}, \texttt{.estimated}) · image loading pipeline (prefetching)}

\end{document}
